\scanpage{54}
% manuscript line 1
\uline{Démonstration} \deleted{En effet} 1) Prouvons p.ex. que $Au$ est
un op. de Hilbert-Schmidt si $u$ \inserted{l'est, et que} \deleted{$\|Au\|_2\unsure{=}$}
$\|Au\|_2\leq\|A\|\,\|u\|_2$. En effet, soit $(e_i)$ une base orthonormale
dans $E$, on a $\deleted{\unsure{\sum\|ue_i\||\sum\|Aue_i\|}}\sum\underset{\unsure{?}}{\|Aue_i\|^2}\leq\deleted{\|A\|}$
$\sum\|A\|^2\|\deleted{\unsure{A}}ue_i\|^2=\|A\|^2\sum\|ue_i\|^2$, d'où \deleted{\unsure{il résulte que}} \inserted{notre assertion.}
L'assertion analogue pour $uA$ en résulte par transposition,
\deleted{si $\sum\|ue_i\|^2$ \unsure{est} fini, il \unsure{en est de même de}}
compte tenu de 2).

% manuscript line 9
\deleted{$\sum$ \illegible{1}} \deleted{2) On a un \unsure{isom.} antilinéaire}
canonique de $\bar E\overset{(2)}{\otimes}F$ sur $\overline{\bar E\overset{(2)}{\otimes}F}$, espace
qui \deleted{peut} s'identifie à $\bar{\bar E}\overset{(2)}{\otimes}\bar F$ (formule 8)
$=E\overset{(2)}{\otimes}\bar F$, \deleted{\unsure{donc}} espace lui-$=$ canoniquemt isomorphe
à $\bar F\overset{(2)}{\otimes}E$. \deleted{\unsure{Il y a ainsi}} En d'autres termes, on a
un \deleted{\unsure{isom.}} isom. d'espaces de Hilbert de
$\deleted{\overline{\bar E\overset{\unsure{(2)}}{\otimes}F}}\ \overline{L^2(E,F)}$ sur $L^2(F,E)$, et on vérifie
aussitôt que \deleted{\unsure{cette}} cette application est donnée
par $\bar u\to u^*$ $(u\in L^2(E,F))$, d'où 2).

% manuscript line 18
3) On sait qu'on a $u=Uh$, où $U$ est
une application partiellt isométrique de $E$ dans $F$,
et $h=Vu$, où $U$ est une application partiellement
isométrique de $F$ dans $E$. On en conclut, \deleted{que} \inserted{en}
\deleted{si \illegible{2} $u$} \inserted{vertu de 1),} que \deleted{\unsure{$h$ et $u$ sont}} \inserted{si l'un des} op. \deleted{de} $u$
et $h$ est de H.-S., il en est de $=$ de
l'autre, \deleted{et \unsure{de norme au plus égale}} \inserted{et a une norme de H.-S. au plus}
égale, d'où en \deleted{fait} \inserted{fait} $\|u\|_2=\|h\|_2$. \deleted{\unsure{On}} D'autre
\deleted{\unsure{est de}} \deleted{\unsure{si $h$ est}} \inserted{pour que $h$}
part, \deleted{pour que} soit un Hilbert-Schmidt, il
faut d'abord que $h$ soit cpt. Alors on sait

% manuscript line 29
% archive sheet number: 55
